\begingroup

\setlength{\LTcapwidth}{\textwidth}

\begin{longtable}{@{}>{\raggedright\arraybackslash}p{7.400cm}>{\raggedleft\arraybackslash}p{1.500cm}>{\raggedleft\arraybackslash}p{1.500cm}>{\raggedleft\arraybackslash}p{2.194cm}@{}}

\caption{Ringkasan penilaian skenario penggunaan CDSS.}

\label{tab:tab6.17_evaluasi_skenario_cdss} \\

\toprule

\textbf{Aspek} & \textbf{Rerata} & \textbf{SB} & \textbf{\% skor $\geq$ 3} \\

\midrule

\endfirsthead



\multicolumn{4}{@{}l}{\small\emph{Tabel \thetable{} (lanjutan)}} \\

\toprule

\textbf{Aspek} & \textbf{Rerata} & \textbf{SB} & \textbf{\% skor $\geq$ 3} \\

\midrule

\endhead



\bottomrule

\endlastfoot



Prediksi dan status kondisi pada CGM & 3,333 & 0,492 & 100,0\% \\



Prediksi finger-stick dan komunikasi keterbatasan & 3,417 & 0,515 & 100,0\% \\



Penolakan prediksi saat data tidak memadai & 3,083 & 0,793 & 91,7\% \\



Menghindari kesan kepastian berlebihan & 3,167 & 0,577 & 91,7\% \\



Menjaga otonomi dokter & 3,333 & 0,492 & 100,0\% \\



Consent dan keamanan data & 3,333 & 0,492 & 100,0\% \\



Dasar rekomendasi cukup jelas & 3,083 & 0,669 & 83,3\% \\



Alur realistis untuk praktik klinis & 3,167 & 0,389 & 100,0\% \\



Karakteristik populasi sumber data memadai & 3,083 & 0,793 & 75,0\% \\

\end{longtable}



\noindent\tabnotestyle Skala 1--4; $n = 12$ responden.\normalsize

\endgroup

